\paragraph{深いレベルからの全射}
主レベルでの空間を \(Y_0\)、一つの正の直線を付けた空間を
\(\widetilde Y_0^0=Y_0\perp\langle1\rangle\) と書く。
深いレベルの法は \(\pi_M^b\) である。
まず実点 \(\omega\) を固定して、\(\phi\) の制限
\(g\mapsto\phi((g\oplus1)\omega)\) が主レベルのすべての形式を
内積で検出することを示す。

各無限素点の二つの最低 \(K\) 型の重さを比較する。
符号反転を含む底本の規約で、共通の正ブロックの重さは零である。
残りの符号は交互で、負側の重さは \([-3\kappa_0,-\kappa_0]\) の中で
狭義に減少し、反対側の重さは正で順序が逆になる。
ここで \(\kappa_0=\kappa(\tau)+\Delta/2\) である。
最低 \(K\) 型と Harish--Chandra パラメータの差は、
各成分から見た他の成分の上下と符号を数えた項であり、
階数以下の補正である。共通の正ブロックではこの数がもとの項を打ち消し、
残りの成分では重さの間隔が補正より大きいので順序が保たれる。

交互の部分について、大きい側の負の重さを \(\mu_i\)、小さい側の
正の重さを \(\mu'_j\) とし、根をそれぞれこれらに持つモニック多項式を
\(P,Q\) とする。実ベクトル
\[
 [F]_i=\frac{F(\mu_i)}{|P'(\mu_i)Q(\mu_i)|^{1/2}}
\]
を使うと、エルミート積は \((-1)^{\deg Q}\) と
\(F_1F_2/(PQ)\) の \(\mu_i\) での留数の和との積である。
\(\deg P=\deg Q+1\) かつ \(\deg Q\) は偶数だから、\([Q]\) のノルムは余次元一の正の直線に
必要な \(1\) になる。\(Q_j=Q/(X-\mu'_j)\) とおく。
\(Q_j/P\) の無限遠での留数は零なので \([Q_j]\) は \([Q]\) に直交する。
\([Q_j]\) と \([Q_{j'}]\) の内積は \(\mu'_j\) で残る留数を取ればよく、
\(j\ne j'\) なら零で、等しい場合は
\((-1)^{\deg Q+1}Q'(\mu'_j)/P(\mu'_j)\) である。
これを絶対値 \(1\) に規格化した列と、共通の正ブロックの恒等写像で
等長包含 \(i_0\) を作る。\((X-\mu'_j)Q_j=Q\) なので、
\(\operatorname{diag}(\mu_i)\) をこの直線の直交補空間へ圧縮すると
\(\operatorname{diag}(\mu'_j)\) になる。
従って \(w=(i_0,[Q])\) として \(w(g\oplus1)w^{-1}\) を使うと、
最低 \(K\) 型の位相の差 \(\operatorname{tr}(\mu(g'-I))-
\operatorname{tr}(\mu'(g-I))\) は零になる。
\leanlinks{\leansrc{New/Determinant/AuxiliaryForms/DeepLevelPoincarePeriod}{最低 K 型の重さの計数}
  \leanlinksep \leansrc{New/Appendix/DiscreteSeries/StationaryPosition}{留数で作る直線と圧縮}}

成分の間隔が \(1\) 以上で絶対値が \(3\kappa_0\) 以下なので、
上の留数ベクトルと等長包含の各成分は \((6\kappa_0)^{M_0+1}\) 以下である。
\(M_0\) は小さい側の交互ブロックの長さである。
従って \(\|w\|^2\le(r+1)^4(6\kappa_0)^{2M_0+2}\le\beta\) とできる。
ここで \(\beta=(\sqrt M)^{\lceil x\rceil}\) と取る。
この位置で二つの離散系列行列係数を対にする。
合計階数の二乗を \(D_0=(2r+1)^2\) とし、
両方の群行列を対角に置いたものを \(J_g\) と書く。
Flensted--Jensen の積分表示では、コンパクト群から離れた部分の
寄与は小行列式の冪の減衰で抑えられる。
コンパクト群の近くでは \(\log(1+z)\) の二次までを取り、
一次位相は上の圧縮で消える。二次のガウス積分が正の主項となる。
位相をそろえた係数の積を \(B(g)\) とし、
\(D(g)=(2/\|J_g\|)^{9D_0}\) とすると、この比較で
\[
 |B(g)|\le D(g),\qquad
 |B(g)|\le 2\RePart B(g)+6\beta^{-50D_0}D(g)
\]
を得る。ここで重さの間隔の余裕は
\(10^{30}(D_0\log\beta)^3\) 以上である。
\(\|J_g-I\|\le(200\beta^5)^{-1}\) では \(|B(g)|\ge1/4\)。
\(\|J_g\|\le\beta+1\) の球は、小球の高々 \(\beta^{16D_0}\) 倍の体積を持つ。
また \(D(g)\) の積分は単位球の体積の \(16^{9D_0}\) 倍以下である。
従って上の誤差を積分しても小球からの寄与を打ち消せず、
\(\displaystyle v\le\beta^{20D_0}|\int B(g)\,dg|\) を得る。
ここで \(v\) は小さい方の群の単位球の体積である。
各無限素点で \(\omega_\tau=w^{-1}\) を取り、これらの積を \(\omega\) とする。
\leanlinks{\leansrc{New/Appendix/DiscreteSeries/GaussianShape}{行列係数のガウス型と誤差}
  \leanlinksep \leansrc{New/Appendix/DiscreteSeries/Periods}{体積比較と積分の下界}
  \leanlinksep \leansrc{New/Appendix/DiscreteSeries/PositionGroup}{余次元一の位置への適用}}

\(Y_0\) の形式 \(f\) がすべての制限に直交し、\(|f|\le S\) とする。
任意のアデール点で \(|f|\le S/2\) を示す。
ここだけで使う深さの尺度を
\(L=\beta^{20r^4(f_E+1)}\) とする。\(f_E=|\Phi|\) である。
最短ベクトルが
\((L^{-2}/(M C_{\mathrm{cusp}}^2))^r\) より短い領域では、
尖点での定数項の消滅により、\(f\) を単位元から \(L^{-2}\) 以内の
実平行移動との差で抑えられる。実再生核の主要部 \(K_\mathrm{maj}\) と
その小移動の差の評価を使うと、この差は
\(2S(\beta/L)\int K_\mathrm{maj}\le S/8\) である。
ここで \(C_{\mathrm{cusp}}\) は階数と不同から定まる格子定数で、
\(\int K_\mathrm{maj}\le L/(16\beta)\) を使った。

最短ベクトルがその閾値以上の点 \(g\) では、\((g\oplus1)\) を中心とし、
\(\widetilde Y_0^0\) の実行列係数を無限部分、
深いレベルの指示関数を有限部分とするポアンカレ級数を取る。
直交の仮定によりその周期は零である。
周期を展開すると、単位剰余類の項はレベルの体積を掛けた
\(f\) の実畳み込みになる。他の剰余類は有理行列の最後の列で区別される。
深い合同条件と最短ベクトルの下界を合わせると、無限素点の最後の列は
互いに \(L^2\) 以上離れる。単位剰余類の列のノルムは \(L\) 以下なので、
他の列の和を \(L^2/2\) から始まる球で数え、実係数の逆冪の減衰を
合わせると \((10/L)^{r^4}\le 10/L\) が残る。
従って実畳み込みは \(S\cdot10/L\) に、平行移動した行列係数の
大きさを表す因子を掛けたもの以下である。
これを \(Y_0\) 自身の再生核でもう一度畳み込む。
Fubini と離散系列の再生性により結果は \(f(g)\) に
先に下から評価した実積分の積を掛けたものである。
その積を割ると、主要核の積分にも含まれる
体積 \(v\) と \(\beta\) の冪が相殺し、
\(|f(g)|\le (10/(16\beta))S\le S/2\) を得る。
全ての点で上界が半分になるので、繰り返して \(f=0\) である。
\leanlinks{\leansrc{New/Determinant/AuxiliaryForms/DeepLevelPoincarePeriod}{最後の列の分離と恒等項}
  \leanlinksep \leansrc{New/Determinant/AuxiliaryForms/DeepLevelVanishing}{尖点の評価と実積分による再生}}

一方、深いレベルを保つ第二の三角形の連鎖は、
第 \ref{sec:exchange-proof} 節の交換と同じ長さの帰納で単射である。
その転置連鎖は像からの制限も含めて全射になる。
実際、転置連鎖の像に直交する形式 \(f\) を上に持ち上げて \(Uf\) とすれば、
随伴の等式で \(\|Uf\|^2=0\) となり、連鎖の単射性から \(f=0\) となる。
この転置連鎖と、直前に構成した \(\omega\) での制限を合成し、
主レベルへ直交射影する。内積で像の直交補空間が零であることを
示したので、得られる写像 \(s\) は主レベルの空間への全射である。
\leanlinks{\leansrc{New/Determinant/AuxiliaryForms/FirstPath}{転置連鎖の全射と制限写像}}
